\chapter{Delta Functions, Quantum Fields, and Conservation Laws}

The simple quantum field introduced earlier is enough to describe scattering, creation, and annihilation in outline, but there is more to extract from it. In particular, it shows how conservation laws enter particle processes. A full course in quantum field theory takes much longer than a few lectures; here the aim is narrower. We shall use the simplest field to see, mechanically, where energy conservation comes from, why a fixed target need not conserve the momentum of the scattered particle, and how phase invariance encodes charge conservation.

Before returning to the field, two pieces of mathematics are needed: the Dirac delta function and the action of creation and annihilation operators on bra vectors.

\section{The Dirac Delta Function}

Call the coordinate \(y\), without deciding whether it will later represent position, wave number, or something else. Begin with an ordinary lump-shaped function concentrated near \(y=a\). Make the lump narrower and narrower. If its height were held fixed, its area would shrink to zero, so raise the height as the width decreases, keeping height times width equal to one. The limiting object is the Dirac delta function \(\delta(y-a)\).

It vanishes away from \(y=a\), is infinitely sharply concentrated at \(y=a\), and has unit area:
\begin{equation}
\delta(y-a)=0\quad (y\neq a),
\qquad
\int_{-\infty}^{\infty}dy\,\delta(y-a)=1.
\end{equation}
The important point is the combination of localization and fixed area. Narrowing the peak does not remove it, because its height increases at the same time.

\begin{figure}[h]
\centering
\includegraphics[width=0.68\textwidth]{lecture_04_figure_02.png}
\caption{Unit-area delta-function spike centered at \(y=a\).}
\lectureframe{00:05:50}
\end{figure}

The same picture can be idealized as a vertical spike:
\begin{figure}[h]
\centering
\begin{tikzpicture}[scale=1.0]
\draw[->] (-3.1,0) -- (3.3,0) node[right] {$y$};
\draw[very thick] (0,0) -- (0,2.9);
\draw (0,-0.15) node[below] {$y=a$};
\node at (1.35,1.65) {$\delta(y-a)$};
\end{tikzpicture}
\caption{Idealized localization of \(\delta(y-a)\) at the zero of its argument.}
\end{figure}

The delta function appears naturally in Fourier integrals. Return for a moment to a periodic spatial interval of total length \(L\). A periodic mode is
\begin{equation}
e^{ikx},
\end{equation}
where \(k\) is one of the allowed wave numbers
\begin{equation}
k_n=\frac{2\pi n}{L},
\qquad n\in\mathbb Z.
\end{equation}
The interval may be represented as \(0\leq x<L\), but it is more symmetric to put the origin in the middle and use
\begin{equation}
-\frac L2\leq x<\frac L2.
\end{equation}
The endpoints are identified: crossing \(L/2\) brings one back at \(-L/2\).

Integrate one allowed mode over the entire periodic interval:
\begin{equation}
I_L(k)=\int_{-L/2}^{L/2}dx\,e^{ikx}.
\end{equation}
If \(k=0\), the integrand is identically one and the answer is \(L\). If \(k\neq0\) is an allowed value, the positive and negative parts of the oscillation cancel over the full number of periods. Thus
\begin{equation}
I_L(k)=
\begin{cases}
L,&k=0,\\[4pt]
0,&k\neq0,
\end{cases}
\qquad
k=\frac{2\pi n}{L}.
\end{equation}

Now regard \(I_L(k)\) as a function on the discrete set of allowed \(k\)'s. The spacing between neighboring values is
\begin{equation}
\Delta k=\frac{2\pi}{L}.
\end{equation}
At \(k=0\) the function has height \(L\); at every other allowed value it is zero. For purposes of drawing the discrete spike, give it one \(k\)-spacing of width. Its area is then
\begin{equation}
(\text{height})(\text{width})
=L\frac{2\pi}{L}=2\pi.
\end{equation}
As \(L\) grows, the height grows, the spacing shrinks, and the area remains \(2\pi\). In the infinite-volume limit the spike becomes \(2\pi\delta(k)\), giving the formal Fourier identity
\begin{equation}
\int_{-\infty}^{\infty}dx\,e^{ikx}=2\pi\delta(k).
\end{equation}
The variables \(x\) and \(k\) occur symmetrically in the exponential. Interchanging their roles gives the companion identity
\begin{equation}
\int_{-\infty}^{\infty}dk\,e^{ikx}=2\pi\delta(x).
\end{equation}
These two elementary rules will soon turn an integral over all interaction times into a condition of energy conservation.

\section{Bras, Kets, and Operators Acting to the Left}

A ket is symbolic notation for a quantum state. For the present purpose, an initial state will be written as a ket and a final state as a bra. Putting a bra and ket together produces an inner product, which is a number. Those numbers, or absolute squares of them, are ultimately what can be compared with experiment.

\begin{classroomqa}
\lecturetimestamp{00:15:30}
\audiencequestion{Had the word ``ket'' actually been introduced before?}
\lecturerresponse{Yes. A ket is one half of the bra-ket notation; the other half, introduced now, is the bra.}
\end{classroomqa}

For the harmonic oscillator, the occupation-number states may be written either as kets \(\lvert n\rangle\) or as bras \(\langle n\rvert\). Their inner product is
\begin{equation}
\langle m\vert n\rangle=\delta_{mn}
=\begin{cases}
1,&m=n,\\
0,&m\neq n.
\end{cases}
\end{equation}
The Kronecker delta \(\delta_{mn}\) is the discrete analogue of the Dirac delta function.

Creation and annihilation operators act on kets in the familiar way:
\begin{align}
a^\dagger\lvert n\rangle&=\sqrt{n+1}\,\lvert n+1\rangle,\\
a\lvert n\rangle&=\sqrt n\,\lvert n-1\rangle.
\end{align}
The creation and annihilation operators may also be denoted \(a^+\) and \(a^-\); here \(a^\dagger\) and \(a\) will be used throughout.

How should the same operators act on bras? The notation puts an operator to the right of a bra when it acts to the left. The definition is chosen so that a matrix element has the same value whether the operator is first applied to the ket or first applied to the bra. Consider
\begin{equation}
\langle n\rvert a^\dagger\lvert m\rangle.
\end{equation}
Acting on the ket gives
\begin{align}
\langle n\rvert a^\dagger\lvert m\rangle
&=\sqrt{m+1}\,\langle n\vert m+1\rangle\\
&=\sqrt{m+1}\,\delta_{n,m+1}.
\end{align}
This is nonzero only if \(n=m+1\), or equivalently if \(m=n-1\). To reproduce the same condition by acting on the bra, \(a^\dagger\) must lower its label:
\begin{equation}
\langle n\rvert a^\dagger
=\sqrt n\,\langle n-1\rvert.
\end{equation}
The numerical factor is \(\sqrt n\), since on the nonzero matrix element \(n=m+1\). The annihilation operator obeys the complementary rule
\begin{equation}
\langle n\rvert a
=\sqrt{n+1}\,\langle n+1\rvert.
\end{equation}
Thus \(a^\dagger\) raises a ket but lowers a bra, while \(a\) lowers a ket but raises a bra.

The advantage is visible in the number-operator matrix element
\begin{equation}
\langle n\rvert a^\dagger a\lvert n\rangle.
\end{equation}
Acting to the right,
\begin{align}
a\lvert n\rangle&=\sqrt n\,\lvert n-1\rangle,\\
a^\dagger\lvert n-1\rangle&=\sqrt n\,\lvert n\rangle,
\end{align}
so
\begin{equation}
\langle n\rvert a^\dagger a\lvert n\rangle
=n\langle n\vert n\rangle=n.
\end{equation}
Acting instead to the left,
\begin{align}
\langle n\rvert a^\dagger&=\sqrt n\,\langle n-1\rvert,\\
\langle n-1\rvert a&=\sqrt n\,\langle n\rvert,
\end{align}
and the same answer follows. One need not decide which direction is more fundamental; the calculus is arranged so that either direction gives the same number.

\section{A Quantum Field in Space and Time}

The reason for this bookkeeping is particle physics. At microscopic scales, the main experimental handle is to collide particles and observe what emerges. An initial state with one or two particles may turn into a state with four, five, seven, nine, or some other number. Quantum fields, built from creation and annihilation operators, provide a language for such processes.

Work first in one spatial dimension. Nothing essential prevents the use of three dimensions: \(x\) and \(k\) simply become vectors with three components. It is also convenient to use units in which
\begin{equation}
\hbar=1.
\end{equation}
Then momentum and wave vector have the same numerical value,
\begin{equation}
\vec p=\vec k.
\end{equation}

\begin{figure}[h]
\centering
\includegraphics[width=0.58\textwidth]{lecture_04_figure_03.png}
\caption{Natural units \(\hbar=1\) and the identification \(\vec p=\vec k\).}
\lectureframe{00:36:56}
\end{figure}

Adopt the convention that \(\psi\) annihilates and \(\psi^\dagger\) creates:
\begin{align}
\psi(x)&=\sum_k a(k)e^{ikx},\\
\psi^\dagger(x)&=\sum_k a^\dagger(k)e^{-ikx}.
\end{align}
The sum is over \(k=2\pi n/L\) in periodic space and becomes continuous in infinite space. The second equation is the Hermitian conjugate of the first. It may also be thought of as the operator analogue of a complex conjugate.\footnote{The blackboard temporarily labels the creation-field expansion \(\sum_k a^\dagger(k)e^{-ikx}\) by \(\Psi\), and later hesitates over the dagger assignment. The labels here are fixed globally so that \(\psi\) annihilates and \(\psi^\dagger\) creates, exactly as in the subsequent scattering calculation.}

Why this definition? At this stage the best answer is that it is useful. It is a simple sum over momentum modes, and its coefficients are precisely the operators that remove or add particles. The detailed justification comes from the dynamics and from the calculations it makes possible.

\editorialnote{A short portion of the source immediately after the conjugate spatial field is introduced is not intelligible enough to support an additional normalization or interpretation statement. The secure content resumes with the usefulness of the definition.}

So far the fields depend on position but not time. Fields oscillate, and each wave number \(k\) has an associated frequency \(\omega(k)\). An object of definite frequency \(\omega\) has time dependence
\begin{equation}
e^{i\omega t}=\cos(\omega t)+i\sin(\omega t).
\end{equation}
Attach the conjugate phases to the two field expansions:
\begin{align}
\psi(x,t)&=\sum_k a(k)e^{ikx}e^{-i\omega(k)t},\\
\psi^\dagger(x,t)&=\sum_k a^\dagger(k)e^{-ikx}e^{i\omega(k)t}.
\end{align}
Now each spatial mode oscillates at the frequency appropriate to its wave number, and the field depends on both space and time.

\section{The Nonrelativistic Field Equation}

The relation between \(\omega\) and \(k\) determines the differential equation obeyed by the field. For a single quantum, the wave-particle relations are
\begin{equation}
E=\hbar\omega,
\qquad
p=\hbar k.
\end{equation}

\begin{figure}[h]
\centering
\includegraphics[width=0.62\textwidth]{lecture_04_figure_04.png}
\caption{Energy-frequency and momentum-wave-number relations \(E=\hbar\omega\) and \(p=\hbar k\).}
\lectureframe{00:41:28}
\end{figure}

With \(\hbar=1\), energy is frequency and momentum is wave number:
\begin{equation}
E=\omega,
\qquad
p=k.
\end{equation}
For a nonrelativistic particle moving much more slowly than light, the kinetic energy and momentum obey
\begin{equation}
E=\frac{p^2}{2m}.
\end{equation}
This is the same relation obtained from \(E=\tfrac12 mv^2\) and \(p=mv\). Combining these relations gives the nonrelativistic dispersion relation
\begin{equation}
\omega(k)=\frac{k^2}{2m}.
\end{equation}

For the creation field, the mode expansion is
\begin{equation}
\psi^\dagger(x,t)=\sum_k a^\dagger(k)e^{-ikx}e^{i\omega(k)t}.
\end{equation}

\begin{figure}[h]
\centering
\includegraphics[width=0.66\textwidth]{lecture_04_figure_05.png}
\caption{Time-dependent mode expansion used to begin the nonrelativistic field equation.}
\lectureframe{00:43:23}
\end{figure}

Differentiate with respect to time:
\begin{equation}
\dot\psi^\dagger(x,t)
=\sum_k i\omega(k)a^\dagger(k)e^{-ikx}e^{i\omega(k)t}.
\end{equation}
Differentiate once with respect to position:
\begin{equation}
\frac{\partial\psi^\dagger}{\partial x}
=\sum_k(-ik)a^\dagger(k)e^{-ikx}e^{i\omega(k)t}.
\end{equation}
There is no direct first-order relation between these derivatives because the time derivative brings down \(\omega(k)\), whereas one spatial derivative brings down \(k\). To see the distinction explicitly, imagine the linear dispersion relation
\begin{equation}
\omega(k)=k.
\end{equation}
The time derivative would then bring down \(i\omega=ik\), while the first spatial derivative brings down \(-ik\). Mode by mode,
\begin{equation}
\frac{\partial\psi^\dagger}{\partial t}
=-\frac{\partial\psi^\dagger}{\partial x}.
\end{equation}
That hypothetical field would satisfy a first-order equation. Here, however, \(\omega=k^2/(2m)\), so a second spatial derivative is required:
\begin{align}
\frac{\partial^2\psi^\dagger}{\partial x^2}
&=\sum_k(-k^2)a^\dagger(k)e^{-ikx}e^{i\omega(k)t}\\
&=-2m\sum_k\omega(k)a^\dagger(k)e^{-ikx}e^{i\omega(k)t}.
\end{align}

\begin{classroomqa}
\lecturetimestamp{00:47:14}
\audiencequestion{Is the minus sign from the spatial derivatives really supposed to be there?}
\lecturerresponse{Yes. Each \(x\)-derivative of \(e^{-ikx}\) brings down \(-ik\), and two derivatives give \((-ik)^2=-k^2\). Substituting \(k^2=2m\omega\) therefore gives \(-2m\omega\).}
\end{classroomqa}

Multiplying the time derivative by \(i\) gives the same mode coefficient as dividing the second spatial derivative by \(2m\):
\begin{align}
i\dot\psi^\dagger
&=-\sum_k\omega(k)a^\dagger(k)e^{-ikx}e^{i\omega(k)t},\\
\frac{1}{2m}\frac{\partial^2\psi^\dagger}{\partial x^2}
&=-\sum_k\omega(k)a^\dagger(k)e^{-ikx}e^{i\omega(k)t}.
\end{align}
Hence
\begin{equation}
i\frac{\partial\psi^\dagger}{\partial t}
=\frac{1}{2m}\frac{\partial^2\psi^\dagger}{\partial x^2},
\end{equation}
or, equivalently,
\begin{equation}
-i\frac{\partial\psi^\dagger}{\partial t}
=-\frac{1}{2m}\frac{\partial^2\psi^\dagger}{\partial x^2}.
\end{equation}

\begin{figure}[h]
\centering
\includegraphics[width=0.72\textwidth]{lecture_04_figure_06.png}
\caption{Time and spatial derivatives of the creation-field mode expansion.}
\lectureframe{00:49:06}
\end{figure}

\begin{classroomqa}
\lecturetimestamp{00:49:43}
\audiencequestion{Is the coefficient \(1/(2m)^2\)?}
\lecturerresponse{No. It is \(1/(2m)\).}
\end{classroomqa}

The conjugate annihilation field obeys the familiar Schr\"odinger form
\begin{equation}
i\frac{\partial\psi}{\partial t}
=-\frac{1}{2m}\frac{\partial^2\psi}{\partial x^2}.
\end{equation}
The two field equations are Hermitian conjugates.\footnote{During the sign check, the board briefly gave the spatial term with the opposite sign. Direct differentiation of the displayed mode phases gives the conjugate pair of equations written here.}

The equation has the same form as the Schr\"odinger equation of elementary quantum mechanics, but the object satisfying it is different. In elementary quantum mechanics, a wavefunction is a complex-valued function whose absolute square is interpreted probabilistically. Here \(\psi\) and \(\psi^\dagger\) are operators: \(\psi\) removes a particle and \(\psi^\dagger\) adds one. The conjugate pair provides a direct link between particles and fields.

When many quanta occupy the relevant modes, a field can behave approximately classically, just as a strongly excited harmonic oscillator can look classical while a one-quantum oscillator does not. Large occupation makes a large classical scale possible and can make appropriate relative fluctuations small.

\editorialnote{Large occupation by itself is not sufficient for a classical field. A phase-coherent semiclassical state is also needed; a number eigenstate has no definite phase and need not have a nonzero mean field.}

\section{A Fixed Target and a Scattering Amplitude}

How are these fields used? Consider the simplest possible scattering process. A particle meets a target so heavy that its recoil will be ignored. The target is fixed in space forever. A particle can be reflected, transmitted, or redirected; in three dimensions its outgoing direction can differ from its incoming direction.

The particle's momentum need not be the same before and after the collision. Momentum is conserved for the complete physical system, but the target absorbs the missing momentum. Once the target's motion is omitted from the mathematical description, the momentum of the particle alone is not conserved.

Its energy behaves differently in the idealized model. Imagine throwing a tennis ball against a wall. A real tennis ball heats and deforms, but if those effects are ignored, it returns with the same kinetic energy even though its momentum reverses. The fixed-target model is designed to reproduce that distinction: the particle's energy is conserved while its momentum may change.

Draw space horizontally and time vertically. The target is the vertical line \(x=0\). An incoming one-particle state of momentum \(k_i\) reaches the target and an outgoing one-particle state of momentum \(k_f\) leaves it.

\begin{figure}[h]
\centering
\begin{tikzpicture}[scale=1.0]
\draw[->] (-3.2,0) -- (3.2,0) node[right] {space};
\draw[->] (0,-0.3) -- (0,4.2) node[above] {time};
\draw[very thick] (0,0) -- (0,3.8);
\node[left] at (-0.15,3.55) {$x=0$};
\draw[thick,->] (-2.2,0.8) -- (0,2.0);
\node[below] at (-1.45,1.05) {$k_i$};
\draw[thick,->] (0,2.0) -- (2.2,3.2);
\node[above] at (1.5,2.9) {$k_f$};
\fill (0,2.0) circle (1.4pt);
\end{tikzpicture}
\caption{A particle of momentum \(k_i\) scatters from a target fixed at \(x=0\) and emerges with momentum \(k_f\).}
\end{figure}

Use the language of absorption and re-emission. The target annihilates the incoming particle and immediately creates the outgoing particle, possibly with a different momentum. This is only a mathematical model of the scattering event; an actual atom, for example, can absorb a photon, remain excited for a time, and emit later.

\begin{classroomqa}
\lecturetimestamp{01:02:28}
\audiencequestion{Why must the particle be absorbed and re-emitted immediately?}
\lecturerresponse{It need not be in nature. Instantaneous absorption and re-emission at one point is the simplifying model being used here.}
\end{classroomqa}

The initial state contains one particle in the \(k_i\) mode and none in any other mode. In occupation-number language every entry is zero except the entry associated with \(k_i\), which is one. It is enough to abbreviate the full list by
\begin{equation}
\lvert k_i\rangle=a^\dagger(k_i)\lvert0\rangle.
\end{equation}
Similarly,
\begin{equation}
\langle k_f\rvert=\langle0\rvert a(k_f).
\end{equation}

Annihilation at the target and time \(t\) is represented by \(\psi(0,t)\); immediate recreation at the same point and time is represented by \(\psi^\dagger(0,t)\). There is no preferred time at which the event occurs. It can happen whenever the incoming particle reaches the target, with the same interaction strength at every time. To express that absence of a special time, integrate over all \(t\):
\begin{equation}
\int_{-\infty}^{\infty}dt\,\psi^\dagger(0,t)\psi(0,t).
\end{equation}
Acting on the initial ket, this operator produces a state. Its overlap with the final bra is the amplitude for the specified transition. The probability is obtained only after taking the absolute square of the amplitude.

One further number is needed. Let \(g\) measure the strength of the interaction between the particle and the target. In this bookkeeping model, write the schematic leading amplitude as\footnote{State-normalization factors and the conventional overall factor of first-order perturbation theory are suppressed; the calculation uses only the operator selection rules and the energy-conserving support.}
\begin{equation}
\mathcal A_{i\to f}
=g\left\langle k_f\left|
\int_{-\infty}^{\infty}dt\,
\psi^\dagger(0,t)\psi(0,t)
\right|k_i\right\rangle.
\end{equation}
The coefficient \(g\) is called a coupling constant. For a charged target scattering a photon, electric charge supplies the relevant measure of interaction strength: increasing the charge increases the tendency to scatter.

Different interactions have very different couplings. A meson reaching a proton undergoes a strong interaction and has a comparatively large scattering strength. A photon scattered by a charged proton has a smaller electromagnetic coupling. A neutrino interacting with a proton has a smaller coupling still.

\begin{classroomqa}
\lecturetimestamp{01:09:09}
\audiencequestion{Is this simple model valid for an electron?}
\lecturerresponse{Only as a simplified stand-in for several kinds of scattering. It should not yet be read as a realistic theory of an electron or of any one particular target.}
\end{classroomqa}

Gravitational scattering of a graviton from a nucleus would be extraordinarily weak. These examples do not share one universal \(g\); particle physics contains many coupling constants.

Nor must an effective coupling always be independent of frequency.

\begin{classroomqa}
\lecturetimestamp{01:10:31}
\audiencequestion{Why take \(g\) to be a constant rather than a function of \(\omega\) or \(k\)?}
\lecturerresponse{It need not always be constant. If a photon is absorbed by an atom, the atom remains excited for a time, and the photon is emitted later, the effective scattering can depend on frequency. The constant \(g\) belongs to the present instantaneous model.}
\end{classroomqa}

For the moment, keep \(g\) constant and calculate the consequences of this simplest choice.

\begin{classroomqa}
\lecturetimestamp{01:12:32}
\audiencequestion{If \(g=1\), must the scattering probability be one?}
\lecturerresponse{No. The expression being calculated is only the leading approximation for small \(g\). If the coupling is not small, higher-order perturbation theory is required; a probability larger than one would signal the breakdown of the one-step approximation, not a physical probability.}
\end{classroomqa}

\section{Collapsing the Mode Sums}

At \(x=0\), the spatial phases equal one, so
\begin{align}
\psi(0,t)&=\sum_k a(k)e^{-i\omega(k)t},\\
\psi^\dagger(0,t)&=\sum_\ell a^\dagger(\ell)e^{i\omega(\ell)t}.
\end{align}

\begin{classroomqa}
\lecturetimestamp{01:13:58}
\audiencequestion{Since the target is at \(x=0\), should the spatial exponential be omitted?}
\lecturerresponse{Yes. At the target \(e^{ikx}=e^{ik0}=1\), and the conjugate spatial factor is also one.}
\end{classroomqa}

Use different dummy labels \(k\) and \(\ell\) in the two sums. Substitution gives
\begin{align}
\mathcal A_{i\to f}
&=g\int_{-\infty}^{\infty}dt
\sum_{\ell,k}
e^{i\omega(\ell)t}e^{-i\omega(k)t}
\langle k_f\rvert a^\dagger(\ell)a(k)\lvert k_i\rangle.
\end{align}
Which term in the \(k\)-sum can survive? An annihilation operator gives zero in every empty mode. The initial state contains a particle only in mode \(k_i\), so \(k=k_i\). Which term in the \(\ell\)-sum can survive? Acting to the left, \(a^\dagger(\ell)\) must remove the particle represented by the final bra, so \(\ell=k_f\). Both sums collapse:
\begin{align}
\mathcal A_{i\to f}
&=g\int_{-\infty}^{\infty}dt\,
e^{i(\omega_f-\omega_i)t}
\langle k_f\rvert a^\dagger(k_f)a(k_i)\lvert k_i\rangle.
\end{align}
The annihilation operator turns the initial one-particle ket into the vacuum,
\begin{equation}
a(k_i)\lvert k_i\rangle=\lvert0\rangle,
\end{equation}
and the creation operator, acting to the left, turns the final one-particle bra into the vacuum bra,
\begin{equation}
\langle k_f\rvert a^\dagger(k_f)=\langle0\rvert.
\end{equation}
Their matrix element is therefore
\begin{equation}
\langle0\vert0\rangle=1.
\end{equation}
All the operator bookkeeping reduces to one; the coupling supplies \(g\), and the nontrivial part left over is the time integral:
\begin{equation}
\mathcal A_{i\to f}
=g\int_{-\infty}^{\infty}dt\,
e^{i(\omega_f-\omega_i)t}.
\end{equation}

\editorialnote{A brief audience question at 01:18:40 is not recoverable from the source. The response distinguishes this one-particle calculation from a relativistic or multiparticle issue; no more precise exchange can be reconstructed securely.}

Now the earlier delta-function identity applies directly:
\begin{equation}
\int_{-\infty}^{\infty}dt\,
e^{i(\omega_f-\omega_i)t}
=2\pi\delta(\omega_f-\omega_i).
\end{equation}
Hence
\begin{equation}
\boxed{
\mathcal A_{i\to f}
=2\pi g\,\delta(\omega_f-\omega_i)
}.
\end{equation}
The amplitude vanishes unless
\begin{equation}
\omega_f=\omega_i.
\end{equation}
Since \(\hbar=1\), this is \(E_f=E_i\). Energy conservation has appeared because every interaction time was treated alike. If the time integral were restricted to a finite interval, it would produce a broadened function rather than an exact delta function. Exact invariance under translations of time is what gives the exact energy-conserving condition.

\begin{classroomqa}
\lecturetimestamp{01:21:30}
\audiencequestion{The amplitude must be squared to obtain a probability. What happens to the delta function when it is squared?}
\lecturerresponse{The squared delta function requires additional care. The point needed here is secure: the result is zero unless the initial and final energies are equal. On the energy-conserving support, the schematic squared prefactor contains \(4\pi^2g^2\).}
\end{classroomqa}

\editorialnote{A normalized transition probability or rate cannot be obtained by simply discarding \(\delta(\omega_f-\omega_i)^2\). It requires finite-time regularization and consistent state normalization. Thus \(4\pi^2g^2\) is retained only as the schematic squared prefactor emphasized in the source.}

Nothing in this point-target interaction distinguishes one final momentum from another except the energy condition. In three dimensions, the model therefore assigns equal leading weight to all outgoing momenta with the same energy---equivalently, to all outgoing directions on the fixed-energy sphere. That isotropy is a property of this very simple point scatterer, not of every real target.

The calculation has supplied three linked ideas. The coefficient \(g\) measures the interaction strength. A bra-operator-ket expression is a scattering amplitude whose absolute square enters the probability. Most importantly, integration over all times produces the energy delta function: time-translation symmetry enforces energy conservation.

\section{What Changes When the Target or Coupling Changes?}

The origin \(x=0\) was chosen only for convenience. Put the target at \(x=x_{\mathrm T}\). The spatial factors in \(\psi^\dagger(x_{\mathrm T},t)\psi(x_{\mathrm T},t)\) give
\begin{equation}
e^{-ik_f x_{\mathrm T}}e^{ik_i x_{\mathrm T}}
=e^{i(k_i-k_f)x_{\mathrm T}}.
\end{equation}
Thus the amplitude becomes
\begin{equation}
\mathcal A_{i\to f}(x_{\mathrm T})
=e^{i(k_i-k_f)x_{\mathrm T}}
\mathcal A_{i\to f}(0).
\end{equation}

\begin{classroomqa}
\lecturetimestamp{01:24:37}
\audiencequestion{Can the position always be redefined so that the target is at \(x=0\)?}
\lecturerresponse{For this single fixed target, yes. If it is left at \(x_{\mathrm T}\neq0\), its position appears as the phase \(e^{i(k_i-k_f)x_{\mathrm T}}\) in the amplitude. Multiplication by the complex-conjugate phase removes it from the probability.}
\end{classroomqa}

Indeed,
\begin{equation}
\left|e^{i(k_i-k_f)x_{\mathrm T}}\right|^2=1.
\end{equation}
The amplitude remembers the target's position, but this one-target probability does not.

Now reverse the asymmetry. Suppose an interaction could occur at any position but only at one distinguished time. An integral over position would give
\begin{equation}
\int_{-\infty}^{\infty}dx\,
e^{i(k_i-k_f)x}
=2\pi\delta(k_i-k_f),
\end{equation}
while fixing the time would remove the energy delta function.

\begin{classroomqa}
\lecturetimestamp{01:26:41}
\audiencequestion{If it did not matter where the interaction happened, but it happened only at one fixed time, would a different quantity be conserved?}
\lecturerresponse{Yes. Integration over all positions would enforce momentum conservation, while the distinguished interaction time would remove exact energy conservation. In a complete two-particle problem, the interaction may occur at any time and any place.}
\end{classroomqa}

This also gives a physical interpretation of explicit time dependence. If the coupling is \(g(t)\), or an outside intervention suddenly makes the scattering possible, not every time is equivalent. The reduced description may then exchange energy with whatever external dynamics were omitted.

\begin{classroomqa}
\lecturetimestamp{01:28:29}
\audiencequestion{Do coupling constants mostly come from experiment, and what does a time-dependent coupling mean?}
\lecturerresponse{Yes, their values are determined experimentally. Explicit time dependence usually means that some moving or changing part of the system has been left out, much as target recoil was left out of the fixed-target model. Energy that appears missing belongs to those omitted degrees of freedom.}
\end{classroomqa}

A moving target is the simplest example.

\begin{classroomqa}
\lecturetimestamp{01:29:23}
\audiencequestion{Is hitting a moving target an example in which the particle's energy is not conserved?}
\lecturerresponse{Yes. If the target's prescribed motion is omitted from the energy bookkeeping, it can supply or remove energy from the particle.}
\end{classroomqa}

The pointlike character of the target is another simplification.

\begin{classroomqa}
\lecturetimestamp{01:29:43}
\audiencequestion{When does the size of the target have to be compared with the size or wavelength of the incident wave?}
\lecturerresponse{In the present model the target has zero size. A finite target size does produce additional effects when compared with the wavelength, but that refinement is postponed.}
\end{classroomqa}

Return now to the moving target. Imagine a wall moving toward a tennis ball. The wall can give energy to the ball. If the wall's motion is prescribed and its energy is omitted, the ball's energy alone is not conserved. Mathematically this is time dependence in the target position rather than necessarily in \(g\).

Creation and annihilation operators have therefore done exactly what was wanted: they describe a transition from one particle state to another, while the spacetime dependence of the fields determines which conservation conditions accompany that transition.

\section{Charge Counting and Phase Invariance}

Each particle species has its own field. Electrons and photons, for example, are not described by the same \(\psi\). To expose charge counting with the algebra already available, temporarily let \(\psi\) stand for an electron-like field: \(\psi\) annihilates one such charged particle and \(\psi^\dagger\) creates one. This is a bosonic stand-in for the bookkeeping; the genuine fermionic electron field requires algebra not yet introduced.

The fixed-target factor
\begin{equation}
\psi^\dagger\psi
\end{equation}
annihilates one particle and creates one. One electron comes in and one goes out, so the charge carried by this species does not change.

Now imagine one electron entering and two leaving from the same point. The corresponding operator product would contain one annihilation field and two creation fields,
\begin{equation}
\psi^\dagger\psi^\dagger\psi.
\end{equation}
For an electron field this would create one net electron and would violate electric charge conservation. Two electrons in and two out can instead be represented by
\begin{equation}
\psi^\dagger\psi^\dagger\psi\psi,
\end{equation}
which has equal numbers of creation and annihilation fields. Two in and three out would require
\begin{equation}
\psi^\dagger\psi^\dagger\psi^\dagger\psi\psi,
\end{equation}
and again changes the net charge. The immediate rule is simple: a charge-conserving interaction contains the same number of \(\psi\)'s and \(\psi^\dagger\)'s.

\begin{classroomqa}
\lecturetimestamp{01:34:27}
\audiencequestion{Does the order of the field operators matter for this charge-conservation test?}
\lecturerresponse{The immediate test concerns the number of \(\psi\) and \(\psi^\dagger\) factors, not their order. Whether differently ordered products are otherwise identical is a separate question.}
\end{classroomqa}

There is a mathematical diagnostic deeper than merely counting. Multiply the field by a constant phase:
\begin{equation}
\psi\longrightarrow e^{i\alpha}\psi.
\end{equation}
Because \(\psi^\dagger\) is the Hermitian conjugate, it transforms with the conjugate phase,
\begin{equation}
\psi^\dagger\longrightarrow e^{-i\alpha}\psi^\dagger.
\end{equation}
An expression with equal numbers of the two fields is unchanged. For example,
\begin{align}
\psi^\dagger\psi
&\longrightarrow
e^{-i\alpha}e^{i\alpha}\psi^\dagger\psi
=\psi^\dagger\psi,\\
\psi^\dagger\psi^\dagger\psi\psi
&\longrightarrow
e^{-2i\alpha}e^{2i\alpha}
\psi^\dagger\psi^\dagger\psi\psi
=\psi^\dagger\psi^\dagger\psi\psi.
\end{align}
An expression with unequal numbers changes. Thus
\begin{equation}
\psi^\dagger\psi^\dagger\psi
\longrightarrow
e^{-i\alpha}\psi^\dagger\psi^\dagger\psi.
\end{equation}
The allowed charge-conserving interactions are precisely the ones invariant under this global change of phase. Phase invariance and charge conservation are two descriptions of the same rule in this simple field language.

The combination
\begin{equation}
\psi+\psi^\dagger
\end{equation}
is not invariant under an arbitrary phase rotation.

\begin{classroomqa}
\lecturetimestamp{01:37:53}
\audiencequestion{Isn't that the real part of \(\psi\)?}
\lecturerresponse{Yes. It is proportional to the real part of \(\psi\), and multiplying \(\psi\) by a phase changes that real part. It is therefore not phase invariant.}
\end{classroomqa}

\begin{classroomqa}
\lecturetimestamp{01:38:05}
\audiencequestion{If an electron strikes a target and many secondary electrons are emitted, how can charge still be conserved?}
\lecturerresponse{The target is not passive in that process. If an atom emits many electrons, the atom itself changes charge and state. Its dynamics must be included, and the target must also be described by quantum fields.}
\end{classroomqa}

At this point three correspondences are visible. Time-translation invariance corresponds to energy conservation. Global phase invariance corresponds to charge conservation. Spatial-translation invariance corresponds to momentum conservation, although the full spatial calculation remains to be carried out.

\section{Operator Order and the Limits of the Lowest-Order Model}

For simplicity, the products above place all creation fields to the left and all annihilation fields to the right. An interleaved product can still have equal numbers and therefore pass the charge-counting test, but operator order can affect the expression in other ways.

\begin{classroomqa}
\lecturetimestamp{01:40:02}
\audiencequestion{Could the same kind of process be written with interleaved factors such as \(\psi\psi^\dagger\psi\psi^\dagger\)?}
\lecturerresponse{It still has equal numbers and so satisfies the charge rule, but it is not generally the same expression as one with all daggers to the left. The rearrangement of different orderings is postponed.}
\end{classroomqa}

\editorialnote{The wording of an audience question at 01:41:10 is too damaged to reconstruct. The recoverable response repeats that equal numbers of \(\psi\) and \(\psi^\dagger\) imply phase invariance and charge conservation, whereas unequal numbers imply neither.}

These links between symmetries and conservation laws are instances of Noether's theorem. In classical mechanics that theorem gives the general framework; here the same connection has been displayed directly through the phases and delta functions of the quantum fields.

The leading fixed-target result also has an obvious domain of validity. A schematic probability prefactor such as \(4\pi^2g^2\) cannot be trusted once it approaches or exceeds unity.

\begin{classroomqa}
\lecturetimestamp{01:43:33}
\audiencequestion{Does the schematic result mean that \(g\) must be smaller than roughly \(1/(2\pi)\)?}
\lecturerresponse{Either \(g\) is small enough for this first term to be useful, or the model is incomplete. A particle can scatter once, return and scatter twice, then three times, and so on. The full amplitude is an infinite series. Only for small \(g\) is the simplest term dominant; otherwise many Feynman graphs must be summed.}
\end{classroomqa}

The simple spacetime pictures are the beginnings of Feynman graphs. Their vertices specify elementary possibilities: one particle in and one out corresponds to a factor such as \(\psi^\dagger\psi\); two in and two out corresponds to a factor such as \(\psi^\dagger\psi^\dagger\psi\psi\). Repeated vertices generate higher powers of the coupling and the higher-order terms in the series.

\begin{classroomqa}
\lecturetimestamp{01:45:39}
\audiencequestion{Can a coupling constant have an imaginary component?}
\lecturerresponse{Complex coupling constants can occur. Irreducible \(T\)-odd phase combinations can signal violation of time-reversal invariance: the time-reversed process need not carry the symmetry-related amplitude.}
\end{classroomqa}

\editorialnote{The QCD \(\theta\) parameter was mentioned in this exchange, but it is a real coefficient multiplying a \(T\)-odd interaction, not an imaginary component of a coupling. A removable phase or an absorptive complex phase does not by itself establish \(T\) violation. The inference applies only to irreducible \(T\)-odd phase combinations.}

\begin{classroomqa}
\lecturetimestamp{01:47:02}
\audiencequestion{Can the outgoing particle be emitted before or after the incoming particle is absorbed, rather than at the same instant?}
\lecturerresponse{That is a fair question, but the instantaneous model does not address it. The more general time ordering is deferred.}
\end{classroomqa}

\begin{classroomqa}
\lecturetimestamp{01:47:44}
\audiencequestion{How is this construction made relativistic?}
\lecturerresponse{The brief indication is to combine the conjugate pieces and include a denominator involving a square root of \(\omega(k)\). The complete construction is postponed.}
\end{classroomqa}

\editorialnote{The brief spoken relativistic formula is not recoverable securely enough to typeset. Only its stated ingredients and deferral are retained.}

Charge conservation must also be distinguished sharply from conservation of rest mass.

\begin{classroomqa}
\lecturetimestamp{01:48:22}
\audiencequestion{Is rest mass conserved in the same structural sense as electric charge?}
\lecturerresponse{No. In Galilean nonrelativistic mechanics, mass is an additive conserved quantity. In relativistic reactions the sum of the individual particles' rest masses need not be conserved. An electron and a positron can annihilate into two photons while electric charge remains conserved.}
\end{classroomqa}

\editorialnote{For a complete closed relativistic system, total four-momentum is conserved, so the invariant mass of that system is fixed. The annihilation example concerns the sum of the individual particles' rest masses, not the invariant mass of the closed system.}

Conservation of additive mass is therefore an effective nonrelativistic fact, not a fundamental rule on the same footing as electric charge conservation. The phase invariance found above expresses the deeper charge symmetry.

The next steps are to replace the bosonic stand-in by fermionic fields and to allow several particle species---electrons, muons, quarks, photons, and others---each with its own field and its own allowed interactions.

\begin{classroomqa}
\lecturetimestamp{01:50:39}
\audiencequestion{Was it legitimate to call the earlier example an electron field when only bosonic fields have been developed so far?}
\lecturerresponse{Only as a bosonic version or stand-in for an electron. The genuine fermionic theory has not yet been introduced.}
\end{classroomqa}
